\documentclass[10pt,a4paper]{article}
\usepackage{amsmath,amssymb,amsthm,graphicx}
% Minimal retained preamble for lane-p2 cards taken from
% sources/ahl-16-fixed-points-2019/editorial-preamble.tex
% (doi:10.5802/ahl.16, Annales Henri Lebesgue 2 (2019) 149-185, CC BY 4.0).
%
% Copied verbatim from the editorial preamble: the package list, the
% theorem declarations under the shared counter and the macro block, so that
% every macro and environment of the retained spans is available.
%
% Omitted and why: the class line (the assembler supplies the class); the
% a4paper geometry line and the fontspec block with its explicit lmroman10
% files (encodings and fonts fight the forced xelatex engine, and the
% assembler supplies the page setup).
\usepackage{amsmath,amssymb,amsthm,mathtools,graphicx,xcolor,hyperref,xurl,xspace,ifthen}
\newtheorem{theo}{Theorem}[section]
\newtheorem{lemm}[theo]{Lemma}
\newtheorem{theorem}[theo]{Theorem}
\newtheorem*{remas*}{Remarks}
\providecommand{\og}{«\nobreakspace}
\providecommand{\fg}{\nobreakspace»}
\newtheorem{lem}[theo]{Lemma}
\newtheorem{lemma}[theo]{Lemma}
\newtheorem{prop}[theo]{Proposition}
\newtheorem{coro}[theo]{Corollary}
\newtheorem{corol}[theo]{Corollary}
\newtheorem{conj}[theo]{Conjecture}
\newtheorem{Proposition}[theo]{Proposition}
\newtheorem{theoreme}[theo]{Theorem}
\theoremstyle{definition}
\newtheorem{defi}[theo]{Definition}
\newtheorem{exam}[theo]{Example}
\newtheorem{exem}[theo]{Example}
\newtheorem{rema}[theo]{Remark}
\newtheorem{nota}[theo]{Notation}
\newtheorem*{rema*}{Remark}
\newtheorem*{defi*}{Definition}
\newtheorem*{theo*}{Theorem}
\newtheorem*{lemm*}{Lemma}
\newtheorem*{prop*}{Proposition}
\newcommand{\enoncename}{}
\newtheorem{corpusenonce}[theo]{\enoncename}
\newenvironment{enonce}[2][]{\renewcommand{\enoncename}{#2}\begin{corpusenonce}}{\end{corpusenonce}}
\newenvironment{enonce*}[2][]{\par\medskip\noindent\textbf{#2.}\itshape}{\par\medskip}
\newenvironment{proofc}[1][\proofname]{\begin{proof}[#1]}{\end{proof}}
\newcommand{\Mc}{\mathcal{M}}
\newcommand{\pointir}{.\ ---\ }
\providecommand{\mathof}[1]{\mathrm{#1}}

\numberwithin{equation}{section}
\numberwithin{figure}{section}
\numberwithin{table}{section}
\def\endappendix{}
%~Mouliné par MaN_auto v.0.21.4 2019-04-30 09:02:09


\usepackage[all]{xy}

\hyphenation{auto-morphism
every
Roelcke
peri-od-ic
peri-od-icy
rep-re-sen-ta-tion
}








\newcommand{\A}{\mathcal A}
\newcommand{\N}{\mathbb{N}}
\newcommand{\Z}{\mathbb{Z}}
\newcommand{\Q}{\mathbb{Q}}
\newcommand{\K}{\mathcal{K}}
\newcommand{\CC}{\mathbb{C}}

\newcommand*{\mA}{\mathbf{A}}
\newcommand*{\mB}{\mathbf{B}}
\newcommand*{\mC}{\mathbf{C}}
\newcommand*{\mD}{\mathbf{D}}
\newcommand*{\mE}{\mathbf{E}}
\newcommand*{\mF}{\mathbf{F}}
\newcommand*{\mH}{\mathbf{H}}
\newcommand*{\mZ}{\mathbf{Z}}
%\newcommand*{\mA}{\mathbf{A}}
%\newcommand*{\mA}{\mathbf{A}}

\newcommand{\restrict}[2]{#1 \upharpoonright #2}

\newcommand{\e}{\varepsilon}

\DeclareMathOperator{\Aut}{Aut}
\DeclareMathOperator{\Stab}{Stab}
\DeclareMathOperator{\Age}{Age}

\newcommand{\actson}{\curvearrowright}

\newcommand{\RUCB}{\mathrm{RUC}_{b}}
\newcommand{\ROB}{\mathrm{Ro}_{b}}
\newcommand{\WAP}{\mathrm{WAP}}
\newcommand{\AP}{\mathrm{AP}}

\graphicspath{{./figures/}}

\newcommand*{\mk}{\mkern -1mu}
\newcommand*{\Mk}{\mkern -2mu}
\newcommand*{\mK}{\mkern 1mu}
\newcommand*{\MK}{\mkern 2mu}

\hypersetup{urlcolor=purple, linkcolor=blue, citecolor=red}

\newcommand*{\romanenumi}{\renewcommand*{\theenumi}{\roman{enumi}}}
\newcommand*{\Romanenumi}{\renewcommand*{\theenumi}{\Roman{enumi}}}
\newcommand*{\alphenumi}{\renewcommand*{\theenumi}{\alph{enumi}}}
\newcommand*{\Alphenumi}{\renewcommand*{\theenumi}{\Alph{enumi}}}
\let\oldexists\exists
\renewcommand*{\exists}{\mathrel{\oldexists}}
\let\oldforall\forall
\renewcommand*{\forall}{\mathrel{\oldforall}}


\usepackage{cleveref}
\crefname{originalproposition}{original proposition}{original propositions}
\Crefname{originalproposition}{Original proposition}{Original propositions}
\makeatletter
\newlabel{prop:Roelcke-precompact}{{4.2}{source}{}{sourcepointer-1}{}}
\newlabel{prop:Roelcke-precompact@cref}{{[originalproposition][0][]4.2}{[1][0][]source}{}{}{}}
\makeatother
\crefname{originalproposition}{original proposition}{original propositions}
\Crefname{originalproposition}{Original proposition}{Original propositions}
\makeatletter
\newlabel{prop:RPExtRoelcke}{{4.9}{source}{}{sourcepointer-2}{}}
\newlabel{prop:RPExtRoelcke@cref}{{[originalproposition][0][]4.9}{[1][0][]source}{}{}{}}
\makeatother
\crefname{originalproposition}{original proposition}{original propositions}
\Crefname{originalproposition}{Original proposition}{Original propositions}
\makeatletter
\newlabel{prop:RoelckeRP}{{4.5}{source}{}{sourcepointer-3}{}}
\newlabel{prop:RoelckeRP@cref}{{[originalproposition][0][]4.5}{[1][0][]source}{}{}{}}
\makeatother
\title{A finite coloring is a Roelcke coloring at a point exactly when one joint embedding of the point determines it}
\author{Lionel Nguyen Van Thé}
\date{doi:10.5802/ahl.16; CC BY 4.0}
\begin{document}
\maketitle
\noindent\textit{Editorial note.} The counted unit is the article's own Proposition prop:Roelcke ext coloring at source lines 567-569 together with its proof at source lines 575-577. The proof block is shared: the next statement, prop:RPExtRoelcke at lines 571-573, is proved by the same block, and both pointers in the proof refer to it. The source is the publisher version of the article in Annales Henri Lebesgue 2 (2019), pages 149-185, DOI 10.5802/ahl.16, released under CC BY 4.0.\par\medskip
\noindent\textit{Editorial note.} The competing candidate of the same block, prop:RPExtRoelcke, is not counted: the article needs both statements on their own, and the counted one is the membership criterion that the proof reads off first, while the adjacent one is the transfer condition quantified over all finite substructures and all joint embeddings. Recording the shared block here keeps the pairing honest, and the adjacent statement is declared as a navigation pointer.\par\medskip
\noindent\textit{Editorial note.} The licence evidence used here is the observation file kept with the source, license-observation.json, which records the licence as CC BY 4.0, points at the inherited publisher-version evidence and states that the landing page displays no licence line; the underlying evidence files are the same ones the earlier cards of this source rely on. No acceptance record is written anywhere.\par\medskip
\noindent\textit{Editorial note.} All three pointers of this card belong to the same chapter of the article and were numbered the way the article itself numbers its statements: the theorem-like environments share one counter attached to the section, so the printed number of each is the section and its place inside that section. The count was made by enumerating every numbered theorem-like environment of the source in order and reading each index off that enumeration: 4.2 for the characterisation of Roelcke precompactness, 4.9 for the extension criterion and 4.5 for the proposition on the Ramsey property of the class. An earlier attempt of this lane used the place in the whole article instead, which is not the printed number, and the cards of this source were all audited against the section-scoped enumeration. The anchors recorded point at each statement's own lines.\par\medskip
\noindent\textit{Editorial note.} The retained preamble is the source's editorial preamble with three omissions only: the class line, which the assembler supplies; the page geometry; and the fontspec block with its explicit font files, which would fight the forced xelatex engine. Everything else, including the theorem declarations under the shared counter and the whole macro block, is retained verbatim.\par\medskip
\noindent\textit{Editorial note.} Neither the retained statement nor the retained proof contains a citation, so this card needs no bibliography entry.\par\medskip
\noindent\textit{Editorial note.} The census package marked this candidate free of missing citation keys and listed the statement span as 567-569 and the proof span as 575-577; both were re-read from the source before the card was written. The proof is dense in the source, three sentences spread over two long lines, and every mathematical move it names is a pointer to another statement of the article, all three of which are declared. Revision r002 of this card replaces the place-in-the-article numbering of the three pointers with the section-scoped printed numbers.\par\medskip
\noindent\textit{Editorial note.} rqual sv-mismatch batch fix 2026-10-14 (Chief ruling: re-assemble unpublished only): source\_version aligned verbatim to the source dir license-observation.json sv (package evidence authority); revision bumped for re-assembly. Recipe from lane-s3.\par\medskip
\noindent\textit{Allowed standard prerequisites (editorial).}\begin{itemize}
\item Roelcke colorings at a point of the Roelcke compactification. For a Fraisse structure with Roelcke precompact automorphism group the article attaches to each point x of the Roelcke compactification the class of colorings that are Roelcke at x, that is, the finite colorings of the copies of a finite substructure that are uniformly continuous at that point of the compactification; this is the class whose Ramsey property the article studies. The author's own definitions are the ones used by the counted statement and proof.. The counted proposition is the criterion that recognises membership in this class combinatorially, by producing a single joint embedding whose pattern determines all the values of the coloring.
\item Joint embeddings and their patterns. For a point x represented by two maps the article considers joint embeddings of the form given by the first map and the second composed with an embedding of a finite structure, and the pattern of such a joint embedding is the function that assigns to every copy the corresponding pair of maps; finite colorings of the copies that only depend on that pattern are exactly the ones the article calls determined by the joint embedding. The author's own definitions are the ones used by the counted statement and proof.. The statement of the counted proposition is phrased with these objects: the coloring belongs to the class at x if and only if some joint embedding of that form determines it.
\end{itemize}\medskip
\ifdefined\hypertarget\hypertarget{sourcepointer-1}{}\fi
\noindent\textit{Source reference.} Original proposition 4.2: sources/ahl-16-fixed-points-2019/body-01.tex, source lines 492--494. Navigation pointer invoked by the retained proof to see that the coloring read off a joint embedding is finite, which is what lets the Ramsey hypothesis apply to it. That proposition is an earlier result of the article and is not retained in this card, so the reference is declared as a pointer carrying the number the article itself prints for that statement. The article declares its theorem-like environments under a single counter attached to the section, so the printed number is the section and the place of the statement inside it, counted by enumerating every numbered theorem-like environment of the source in order: it is 4.2; the cited statement is an allowed input and none of it is reproduced inside the span.\par\medskip
\ifdefined\hypertarget\hypertarget{sourcepointer-2}{}\fi
\noindent\textit{Source reference.} Original proposition 4.9: sources/ahl-16-fixed-points-2019/body-01.tex, source lines 571--573. Navigation pointer invoked by the retained proof for the extension criterion, which shows that the coloring read off a joint embedding belongs to the class at the point and which also supplies the converse direction. The adjacent statement shares this proof block but is not retained in the card, so it is declared as a pointer carrying the number the article itself prints for that statement, that is the section and the place of the statement inside it, counted by enumerating every numbered theorem-like environment of the source in order: it is 4.9; the cited statement is an allowed input and none of it is reproduced inside the span.\par\medskip
\ifdefined\hypertarget\hypertarget{sourcepointer-3}{}\fi
\noindent\textit{Source reference.} Original proposition 4.5: sources/ahl-16-fixed-points-2019/body-01.tex, source lines 519--521. Navigation pointer invoked by the retained proof for the converse direction, whose argument captures a finite set of colorings of the class by one single such coloring. That proposition is an earlier result of the article and is not retained in this card, so the reference is declared as a pointer carrying the number the article itself prints for that statement, that is the section and the place of the statement inside it, counted by enumerating every numbered theorem-like environment of the source in order: it is 4.5; the cited statement is an allowed input and none of it is reproduced inside the span.\par\medskip
\begin{prop}\label{prop:Roelcke ext coloring}
Let $\mF$ be a Fra\"{\i}ss\'e structure with Roelcke-precompact automorphism group, $x=[\phi_{1},\phi_{2}]\in R(\Aut(\mF))$, $\mA\in \Age(\mF)$ and $\chi$ be a finite coloring of $\binom{\mF}{\mA}$. Then $\chi\in \ROB(\Aut(\mF))_{x}$ iff there is $\mZ\in \Age(\mF)$ and a joint embedding of the form $\langle \phi_{1}, \phi_{2}\circ z \rangle$ of $\mF$ and $\mZ$ such that $\chi(a)$ depends only on the joint embedding pattern $[\phi_{1}\circ a,\phi_{2}\circ z]$.
\end{prop}

\begin{proof}
Assume that $\mF$ has the Ramsey property for colorings in $\ROB(\Aut(\mF))_{x}$, and fix $\mA, \mB, \mZ\in \Age(\mF)$ together with a joint embedding of $\mF$ and $\mZ$ of the form $\langle \phi_{1}, \phi_{2}\circ z\rangle$. Then the coloring defined on $\binom{\mF}{\mA}$ by $a\mapsto [\phi_{1} \circ a,\phi_{2}\circ z]$ is finite by Proposition~\ref{prop:Roelcke-precompact}, and is in $\ROB(\Aut(\mF))_{x}$ by Proposition~\ref{prop:Roelcke ext coloring}. The conclusion follows. The converse is an immediate consequence of Proposition~\ref{prop:Roelcke ext coloring}, and of the fact that any finite set $\mathcal C$ of finite colorings in $\ROB(\Aut(\mF))_{x}$ can be captured by one single such coloring, as in the proof of Proposition~\ref{prop:RoelckeRP}.
\end{proof}

\end{document}
